\RequirePackage{plautopatch}
\documentclass[platex,dvipdfmx,a4paper]{jsarticle}

\usepackage{amsmath}
\usepackage{amsfonts}
\usepackage{amssymb}

\usepackage[dvipdfmx]{graphicx}
\usepackage{float}

\usepackage{tikz}
\usepackage{caption}

% 「%」は以降の内容を「改行コードも含めて」無視するコマンド
\usepackage[%
  dvipdfmx,% 欧文ではコメントアウトする
  setpagesize=false,%
  bookmarks=true,%
  bookmarksdepth=tocdepth,%
  bookmarksnumbered=true,%
  colorlinks=false,%
  pdftitle={},%
  pdfsubject={},%
  pdfauthor={},%
  pdfkeywords={}%
]{hyperref}
% PDFのしおり機能の日本語文字化けを防ぐ((u)pLaTeXのときのみかく)
\usepackage{pxjahyper}

\usepackage{framed}
\setlength{\OuterFrameSep}{0pt}

\newcommand{\hb}[2]{\frac{\partial #1}{\partial #2}}
\newcommand{\bbn}[2]{\frac{d #1}{d #2}}
\newcommand{\nhb}[3]{\frac{\partial^{#1} #2}{\partial #3^{#1}}}
\newcommand{\nbbn}[3]{\frac{d^{#1} #2}{d #3^{#1}}}
\newcommand{\vectxy}[2]{\begin{bmatrix}#1\\#2\end{bmatrix}}
\newcommand{\vectxyz}[3]{\begin{bmatrix}#1\\#2\\#3\end{bmatrix}}
\newcommand{\vect}[1]{\boldsymbol{\mathrm{#1}}}

\newcommand{\ssum}[3]{\sum_{#1=#2}^#3}

\newcommand{\veff}{V_{\mathrm{eff}}(r)}

\hypersetup{colorlinks=true}
\hypersetup{linkcolor=red}

\title{物理数学II期末レポート}
\author{25B00729　嶋田 悠音}

\begin{document}
\maketitle

  \subsection*{問題1}

  この問題では、周期$2\pi$の関数
  \begin{align*}
    f(x)=|x|\quad(-\pi\leq x<\pi)
  \end{align*}
  のFourier級数展開を求める。

  まず、$f(x)=|x|$は偶関数であるため：
  \begin{align*}
    f(-x)=f(x)
  \end{align*}
  が成り立つ。したがって、一般のFourier級数
  \begin{align*}
    f(x)=\frac{a_0}{2}+\sum_{n=1}^{\infty}a_n\cos(nx)+\sum_{n=1}^{\infty}b_n\sin(nx)
  \end{align*}
  において、正弦項の係数はすべてゼロとなる：
  \begin{align*}
    b_n
    &=\frac{1}{\pi}\int_{-\pi}^{\pi}|x|\sin(nx)\,dx\\
    &=0。
  \end{align*}

  次に、定数項の係数$a_0$を求める。被積分関数が偶関数であることを用いる：
  \begin{align*}
    a_0
    &=\frac{1}{\pi}\int_{-\pi}^{\pi}|x|\,dx\\
    &=\frac{2}{\pi}\int_0^\pi x\,dx\\
    &=\frac{2}{\pi}\left[\frac{x^2}{2}\right]_0^\pi\\
    &=\pi。
  \end{align*}
  したがって、Fourier級数の定数項は次のようになる：
  \begin{align*}
    \frac{a_0}{2}=\frac{\pi}{2}。
  \end{align*}

  続いて、$n\geq1$に対する余弦項の係数$a_n$を求める。被積分関数が偶関数であることを用いる：
  \begin{align*}
    a_n
    &=\frac{1}{\pi}\int_{-\pi}^{\pi}|x|\cos(nx)\,dx\\
    &=\frac{2}{\pi}\int_0^\pi x\cos(nx)\,dx。
  \end{align*}
  ここで、部分積分を用いる：
  \begin{align*}
    a_n
    &=\frac{2}{\pi}\left\{
      \left[\frac{x\sin(nx)}{n}\right]_0^\pi
      -\frac{1}{n}\int_0^\pi\sin(nx)\,dx
    \right\}\\
    &=\frac{2}{\pi}\left\{
      \left[\frac{x\sin(nx)}{n}\right]_0^\pi
      +\left[\frac{\cos(nx)}{n^2}\right]_0^\pi
    \right\}\\
    &=\frac{2}{\pi n^2}\left\{\cos(n\pi)-1\right\}\\
    &=\frac{2}{\pi n^2}\left\{(-1)^n-1\right\}。
  \end{align*}
  よって、$n$が偶数の場合には$a_n=0$であり、$n$が奇数の場合には次の式が成り立つ：
  \begin{align*}
    a_n=-\frac{4}{\pi n^2}\quad(n\text{：奇数})。
  \end{align*}

  以上のFourier係数を代入することにより、求めるFourier級数を得る：
  \begin{align*}
    |x|
    =\frac{\pi}{2}
    -\frac{4}{\pi}\sum_{m=0}^{\infty}
    \frac{\cos\{(2m+1)x\}}{(2m+1)^2}
    \quad(-\pi\leq x<\pi)。
  \end{align*}
  関数$|x|$を周期$2\pi$で延長した関数は連続であるため、このFourier級数はすべての$x$において元の周期関数へ収束する。以上より、$f(x)=|x|$のFourier級数展開が求められた。

  \subsection*{問題2}

  区間$-\infty<x<\infty$および$-\infty<t<\infty$において、熱方程式のGreen関数が満たす方程式
  \begin{align*}
    a^2\hb{^2G(x,t;x_0,t_0)}{x^2}
    -\hb{G(x,t;x_0,t_0)}{t}
    =-\delta(x-x_0)\delta(t-t_0)
  \end{align*}
  を考える。ただし、$a>0$とする。

  \subsubsection*{(1)}

  この小問では、$x$および$t$についてFourier変換を行い、Fourier空間におけるGreen関数$g(k,\omega)$を求める。

  Green関数とデルタ関数を次のようにFourier変換する：
  \begin{align*}
    G(x,t;x_0,t_0)
    &=\frac{1}{(2\pi)^2}
    \int_{-\infty}^{\infty}dk
    \int_{-\infty}^{\infty}d\omega\,
    g(k,\omega)e^{ik(x-x_0)-i\omega(t-t_0)},\\
    \delta(x-x_0)\delta(t-t_0)
    &=\frac{1}{(2\pi)^2}
    \int_{-\infty}^{\infty}dk
    \int_{-\infty}^{\infty}d\omega\,
    e^{ik(x-x_0)-i\omega(t-t_0)}。
  \end{align*}

  Green関数への$x$および$t$による微分は、$x,t$が積分に関係しないことからそれぞれ被積分関数に作用する：
  \begin{align*}
    \hb{^2}{x^2}
    e^{ik(x-x_0)-i\omega(t-t_0)}
    &=-k^2e^{ik(x-x_0)-i\omega(t-t_0)},\\
    \hb{}{t}
    e^{ik(x-x_0)-i\omega(t-t_0)}
    &=-i\omega e^{ik(x-x_0)-i\omega(t-t_0)}。
  \end{align*}
  これらを用いてFourier変換されたGreen関数を与えられた方程式へ代入し、両辺のFourier係数を比較する：
  \begin{align*}
    \left(-a^2k^2+i\omega\right)g(k,\omega)=-1。
  \end{align*}
  したがって、Fourier空間におけるGreen関数は次のようになる：
  \begin{align*}
    g(k,\omega)
    &=\frac{1}{a^2k^2-i\omega}\\
    &=\frac{i}{\omega+ia^2k^2}。
  \end{align*}
  以上より、$g(k,\omega)$が求められた。

  \subsubsection*{(2)}

  この小問では、時間方向の逆Fourier変換
  \begin{align*}
    \int_{-\infty}^{\infty}\frac{d\omega}{2\pi}
    g(k,\omega)e^{-i\omega(t-t_0)}
  \end{align*}
  を計算する。

  時間差を次のように置く：
  \begin{align*}
    \tau=t-t_0。
  \end{align*}
  (1)で得た$g(k,\omega)$を代入する：
  \begin{align*}
    I(k,\tau)
    &=\int_{-\infty}^{\infty}\frac{d\omega}{2\pi}
    \frac{i\,e^{-i\omega\tau}}{\omega+ia^2k^2}。
  \end{align*}
  被積分関数は複素$\omega$平面の下半平面に単純極
  \begin{align*}
    \omega=-ia^2k^2
  \end{align*}
  をもつ。

  まず、$\tau>0$の場合を考える。このとき、指数因子$e^{-i\omega\tau}$が減衰する下半平面に積分路を閉じる。閉曲線の向きは時計回りであるため、留数定理より：
  \begin{align*}
    I(k,\tau)
    &=-i\mathop{\mathrm{Res}}_{\omega=-ia^2k^2}
    \frac{i\,e^{-i\omega\tau}}{\omega+ia^2k^2}\\
    &=-i\left\{i\,e^{-a^2k^2\tau}\right\}\\
    &=e^{-a^2k^2\tau}。
  \end{align*}
  一方、$\tau<0$の場合には上半平面に積分路を閉じるが、上半平面には極が存在しないため、積分はゼロとなる。したがって、両者をHeaviside関数によってまとめる：
  \begin{align*}
    \int_{-\infty}^{\infty}\frac{d\omega}{2\pi}
    g(k,\omega)e^{-i\omega(t-t_0)}
    =
    \theta(t-t_0)e^{-a^2k^2(t-t_0)}。
  \end{align*}
  以上より、時間方向の逆Fourier変換が求められた。

  \subsubsection*{(3)}

  この小問では、(2)の結果を用いて実空間におけるGreen関数を求める。

  (2)の結果をGreen関数のFourier表示へ代入する：
  \begin{align*}
    G(x,t;x_0,t_0)
    =
    \theta(t-t_0)
    \int_{-\infty}^{\infty}\frac{dk}{2\pi}
    e^{-a^2k^2(t-t_0)+ik(x-x_0)}。
  \end{align*}
  ここで、再び
  \begin{align*}
    \tau=t-t_0,
    \qquad
    X=x-x_0
  \end{align*}
  と置く。$\tau>0$に対して、指数部分を平方完成する：
  \begin{align*}
    -a^2\tau k^2+ikX
    =-a^2\tau\left(k-\frac{iX}{2a^2\tau}\right)^2
    -\frac{X^2}{4a^2\tau}。
  \end{align*}
  したがって：
  \begin{align*}
    G(x,t;x_0,t_0)
    &=
    \frac{\theta(\tau)}{2\pi}
    e^{-\frac{X^2}{4a^2\tau}}
    \int_{-\infty}^{\infty}
    \exp\left\{
      -a^2\tau\left(k-\frac{iX}{2a^2\tau}\right)^2
    \right\}dk。
  \end{align*}
  被積分関数は整関数であり、十分遠方で減衰するため、積分路を実軸へ平行移動できる。Gauss積分
  \begin{align*}
    \int_{-\infty}^{\infty}e^{-\alpha k^2}\,dk
    =\sqrt{\frac{\pi}{\alpha}}
    \qquad(\alpha>0)
  \end{align*}
  を$\alpha=a^2\tau$として用いる：
  \begin{align*}
    \int_{-\infty}^{\infty}e^{-a^2\tau k^2}\,dk
    =\frac{\sqrt{\pi}}{a\sqrt{\tau}}。
  \end{align*}
  よって、Green関数は次のようになる：
  \begin{align*}
    G(x,t;x_0,t_0)
    =
    \frac{\theta(t-t_0)}
    {2a\sqrt{\pi(t-t_0)}}
    \exp\left\{
      -\frac{(x-x_0)^2}{4a^2(t-t_0)}
    \right\}。
  \end{align*}

  \subsection*{問題3}
  
  Laguerre多項式$L_n(x)$の生成関数
  \begin{align*}
    G(x,t)
    =\frac{1}{1-t}\exp\left(-\frac{xt}{1-t}\right)
    =\sum_{n=0}^{\infty}L_n(x)t^n
  \end{align*}
  を用いて、以下の各問を解く。

  \subsubsection*{(1)}

  この小問では、Laguerre多項式の生成関数を用いて、$L_0(x)$および$L_1(x)$を求める。

  まず、生成関数に$t=0$を代入する：
  \begin{align*}
    G(x,0)
    &=\exp(0)\\
    &=1。
  \end{align*}
  一方、生成関数の級数表示に$t=0$を代入すると：
  \begin{align*}
    G(x,0)=L_0(x)。
  \end{align*}
  したがって、次の式を得る：
  \begin{align*}
    L_0(x)=1。
  \end{align*}

  次に、$L_1(x)$を求める。生成関数の級数表示を$t$で偏微分する：
  \begin{align*}
    \hb{G(x,t)}{t}
    =\sum_{n=1}^{\infty}nL_n(x)t^{n-1}。
  \end{align*}
  ここに$t=0$を代入すると、$n=1$の項だけが残るため：
  \begin{align*}
    \left.\hb{G(x,t)}{t}\right|_{t=0}=L_1(x)。
  \end{align*}

  そこで、生成関数の具体的な表式を$t$で偏微分する。積の微分および合成関数の微分を用いる：
  \begin{align*}
    \hb{G(x,t)}{t}
    &=
    \hb{}{t}\left\{
      \frac{1}{1-t}\exp\left(-\frac{xt}{1-t}\right)
    \right\}\\
    &=
    \frac{1}{(1-t)^2}\exp\left(-\frac{xt}{1-t}\right)
    +\frac{1}{1-t}\exp\left(-\frac{xt}{1-t}\right)
    \hb{}{t}\left(-\frac{xt}{1-t}\right)。
  \end{align*}
  指数関数の指数部分を偏微分する。商の微分を用いると：
  \begin{align*}
    \hb{}{t}\left(-\frac{xt}{1-t}\right)
    &=-x\hb{}{t}\left(\frac{t}{1-t}\right)\\
    &=-x\frac{(1-t)+t}{(1-t)^2}\\
    &=-\frac{x}{(1-t)^2}。
  \end{align*}
  したがって、生成関数の偏微分は次のようになる：
  \begin{align*}
    \hb{G(x,t)}{t}
    &=
    \frac{1}{(1-t)^2}\exp\left(-\frac{xt}{1-t}\right)
    -\frac{x}{(1-t)^3}\exp\left(-\frac{xt}{1-t}\right)\\
    &=
    \exp\left(-\frac{xt}{1-t}\right)
    \left\{\frac{1}{(1-t)^2}-\frac{x}{(1-t)^3}\right\}。
  \end{align*}
  ここに$t=0$を代入する：
  \begin{align*}
    L_1(x)
    &=\left.\hb{G(x,t)}{t}\right|_{t=0}\\
    &=\exp(0)(1-x)\\
    &=1-x。
  \end{align*}
  以上より、求めるLaguerre多項式は次のようになる：
  \begin{align*}
    L_0(x)=1,
    \qquad
    L_1(x)=1-x。
  \end{align*}

  を得る。以上より、$L_0(x)$および$L_1(x)$が求められた。

  \subsubsection*{(2)}

  この小問では、生成関数の積を含む積分
  \begin{align*}
    \int_0^\infty e^{-x}G(x,t)G(x,s)\,dx
  \end{align*}
  を計算する。

  生成関数の具体的な表式を代入する：
  \begin{align*}
    &\int_0^\infty e^{-x}G(x,t)G(x,s)\,dx\\
    &\quad=
    \frac{1}{(1-t)(1-s)}
    \int_0^\infty
    \exp\left\{
      -x-\frac{xt}{1-t}-\frac{xs}{1-s}
    \right\}dx。
  \end{align*}
  指数部分の$x$の係数を整理する：
  \begin{align*}
    1+\frac{t}{1-t}+\frac{s}{1-s}
    &=
    \frac{(1-t)(1-s)+t(1-s)+s(1-t)}
    {(1-t)(1-s)}\\
    &=
    \frac{1-st}{(1-t)(1-s)}。
  \end{align*}
  したがって：
  \begin{align*}
    &\int_0^\infty e^{-x}G(x,t)G(x,s)\,dx\\
    &\quad=
    \frac{1}{(1-t)(1-s)}
    \int_0^\infty
    \exp\left\{
      -\frac{1-st}{(1-t)(1-s)}x
    \right\}dx。
  \end{align*}
  $|s|<1$および$|t|<1$の範囲では積分が収束する。指数関数の積分を実行する：
  \begin{align*}
    \int_0^\infty e^{-x}G(x,t)G(x,s)\,dx
    &=
    \frac{1}{(1-t)(1-s)}
    \frac{(1-t)(1-s)}{1-st}\\
    &=
    \frac{1}{1-st}。
  \end{align*}
  以上より、求める積分は次のようになる：
  \begin{align*}
    \int_0^\infty e^{-x}G(x,t)G(x,s)\,dx
    =\frac{1}{1-st}。
  \end{align*}

  \subsubsection*{(3)}

  この小問では、(2)の結果を用いて、Laguerre多項式の内積
  \begin{align*}
    \int_0^\infty e^{-x}L_n(x)L_m(x)\,dx
  \end{align*}
  を評価する。

  生成関数の級数表示を用いる：
  \begin{align*}
    G(x,t)=\sum_{n=0}^{\infty}L_n(x)t^n,
    \qquad
    G(x,s)=\sum_{m=0}^{\infty}L_m(x)s^m。
  \end{align*}
  これらを積分へ代入する：
  \begin{align*}
    \int_0^\infty e^{-x}G(x,t)G(x,s)\,dx
    &=
    \int_0^\infty e^{-x}
    \left\{\sum_{n=0}^{\infty}L_n(x)t^n\right\}
    \left\{\sum_{m=0}^{\infty}L_m(x)s^m\right\}dx\\
    &=
    \sum_{n=0}^{\infty}\sum_{m=0}^{\infty}
    t^ns^m
    \int_0^\infty e^{-x}L_n(x)L_m(x)\,dx。
  \end{align*}
  一方、(2)の結果を$st$の関数と見てテイラー展開する：
  \begin{align*}
    \frac{1}{1-st}
    =\sum_{\ell=0}^{\infty}(st)^\ell
    =\sum_{\ell=0}^{\infty}s^\ell t^\ell。
  \end{align*}
  したがって：
  \begin{align*}
    \sum_{n=0}^{\infty}\sum_{m=0}^{\infty}
    t^ns^m
    \int_0^\infty e^{-x}L_n(x)L_m(x)\,dx
    =
    \sum_{\ell=0}^{\infty}s^\ell t^\ell。
  \end{align*}
  $s^mt^n$の係数を比較する。このとき、右辺では$m=n$の場合にのみ係数が$1$となり、$m\neq n$の場合には係数がゼロとなる。したがって：
  \begin{align*}
    \int_0^\infty e^{-x}L_n(x)L_m(x)\,dx
    =\delta_{nm}。
  \end{align*}
  以上より、Laguerre多項式は重み関数$e^{-x}$に関して区間$[0,\infty)$上で直交し、本問の規格化では規格化定数が$1$となることが示された。

  \subsection*{問題4}

  整数$n=0,\pm1,\pm2,\dots$に対して、Bessel関数$J_n(z)$が次の積分表示によって定義されているとする：
  \begin{align*}
    J_n(z)
    =\frac{1}{\pi}\int_0^\pi\cos(z\sin\theta-n\theta)\,d\theta。
  \end{align*}
  以下では、この積分表示からBessel関数の対称性、漸化式および微分方程式を導く。

  \subsubsection*{(1)}

  この小問では、整数$n$に対して次の関係を示す：
  \begin{align*}
    J_{-n}(z)=(-1)^nJ_n(z)。
  \end{align*}

  積分表示において$n$を$-n$へ置き換える：
  \begin{align*}
    J_{-n}(z)
    =\frac{1}{\pi}\int_0^\pi\cos(z\sin\theta+n\theta)\,d\theta。
  \end{align*}
  ここで、積分変数を
  \begin{align*}
    \theta=\pi-\varphi
  \end{align*}
  と変換する。このとき、$d\theta=-d\varphi$であり：
  \begin{align*}
    \sin\theta=\sin(\pi-\varphi)=\sin\varphi
  \end{align*}
  が成り立つ。したがって：
  \begin{align*}
    J_{-n}(z)
    &=\frac{1}{\pi}\int_0^\pi
    \cos\{z\sin\varphi+n(\pi-\varphi)\}\,d\varphi\\
    &=\frac{1}{\pi}\int_0^\pi
    \cos\{z\sin\varphi-n\varphi+n\pi\}\,d\varphi。
  \end{align*}
  $n$は整数であるため、$\sin(n\pi)=0$および$\cos(n\pi)=(-1)^n$が成り立つ。加法定理を用いる：
  \begin{align*}
    \cos\{z\sin\varphi-n\varphi+n\pi\}
    =(-1)^n\cos(z\sin\varphi-n\varphi)。
  \end{align*}
  よって：
  \begin{align*}
    J_{-n}(z)
    &=\frac{(-1)^n}{\pi}
    \int_0^\pi\cos(z\sin\varphi-n\varphi)\,d\varphi\\
    &=(-1)^nJ_n(z)。
  \end{align*}
  以上より、$J_{-n}(z)=(-1)^nJ_n(z)$が示された。

  \subsubsection*{(2)}

  この小問では、Bessel関数の導関数について次の漸化式を示す：
  \begin{align*}
    \bbn{J_n(z)}{z}
    =\frac{1}{2}\left\{J_{n-1}(z)-J_{n+1}(z)\right\}。
  \end{align*}

  Bessel関数の積分表示を$z$で微分する。積分区間は有限であり、被積分関数とその$z$微分は連続であるため、積分と微分を交換できる：
  \begin{align*}
    \bbn{J_n(z)}{z}
    &=-\frac{1}{\pi}\int_0^\pi
    \sin(z\sin\theta-n\theta)\sin\theta\,d\theta。
  \end{align*}

  一方、$J_{n-1}(z)$と$J_{n+1}(z)$の差を積分表示によって書く。ここで、
  \begin{align*}
    A=z\sin\theta-n\theta
  \end{align*}
  と置くと：
  \begin{align*}
    J_{n-1}(z)-J_{n+1}(z)
    &=\frac{1}{\pi}\int_0^\pi
    \left\{\cos(A+\theta)-\cos(A-\theta)\right\}d\theta。
  \end{align*}
  三角関数の加法定理より：
  \begin{align*}
    \cos(A+\theta)-\cos(A-\theta)
    =-2\sin A\sin\theta
  \end{align*}
  が成り立つ。したがって：
  \begin{align*}
    J_{n-1}(z)-J_{n+1}(z)
    &=-\frac{2}{\pi}\int_0^\pi
    \sin(z\sin\theta-n\theta)\sin\theta\,d\theta\\
    &=2\bbn{J_n(z)}{z}。
  \end{align*}
  よって：
  \begin{align*}
    \bbn{J_n(z)}{z}
    =\frac{1}{2}\left\{J_{n-1}(z)-J_{n+1}(z)\right\}。
  \end{align*}
  以上より、導関数に関する漸化式が示された。

  \subsubsection*{(3)}

  この小問では、次の漸化式を示す：
  \begin{align*}
    \frac{n}{z}J_n(z)
    =\frac{1}{2}\left\{J_{n-1}(z)+J_{n+1}(z)\right\}。
  \end{align*}

  まず、位相
  \begin{align*}
    A(\theta)=z\sin\theta-n\theta
  \end{align*}
  を導入する。$\sin A(\theta)$を$\theta$で微分する：
  \begin{align*}
    \bbn{}{\theta}\sin A(\theta)
    =\left(z\cos\theta-n\right)\cos A(\theta)。
  \end{align*}
  この式を$0$から$\pi$まで積分する：
  \begin{align*}
    \int_0^\pi
    \left(z\cos\theta-n\right)
    \cos(z\sin\theta-n\theta)\,d\theta
    =
    \left[\sin(z\sin\theta-n\theta)\right]_0^\pi。
  \end{align*}
  $n$が整数であるため、右辺は次のようにゼロとなる：
  \begin{align*}
    \left[\sin(z\sin\theta-n\theta)\right]_0^\pi
    =\sin(-n\pi)-\sin0
    =0。
  \end{align*}
  したがって：
  \begin{align*}
    n\int_0^\pi\cos(z\sin\theta-n\theta)\,d\theta
    =
    z\int_0^\pi
    \cos\theta\cos(z\sin\theta-n\theta)\,d\theta。
  \end{align*}
  両辺を$\pi z$で割る：
  \begin{align*}
    \frac{n}{z}J_n(z)
    =
    \frac{1}{\pi}\int_0^\pi
    \cos\theta\cos(z\sin\theta-n\theta)\,d\theta。
  \end{align*}

  一方、$A=z\sin\theta-n\theta$と置き、加法定理
  \begin{align*}
    \cos(A+\theta)+\cos(A-\theta)=2\cos A\cos\theta
  \end{align*}
  を用いる：
  \begin{align*}
    J_{n-1}(z)+J_{n+1}(z)
    &=\frac{1}{\pi}\int_0^\pi
    \left\{\cos(A+\theta)+\cos(A-\theta)\right\}d\theta\\
    &=\frac{2}{\pi}\int_0^\pi
    \cos A\cos\theta\,d\theta\\
    &=\frac{2n}{z}J_n(z)。
  \end{align*}
  したがって：
  \begin{align*}
    \frac{n}{z}J_n(z)
    =\frac{1}{2}\left\{J_{n-1}(z)+J_{n+1}(z)\right\}。
  \end{align*}
  以上より、Bessel関数の第2の漸化式が示された。

  \subsubsection*{(4)}

  この小問では、(2)および(3)で得た漸化式を用いて、$J_n(z)$がBesselの微分方程式
  \begin{align*}
    \nbbn{2}{J_n}{z}
    +\frac{1}{z}\bbn{J_n}{z}
    +\left(1-\frac{n^2}{z^2}\right)J_n
    =0
  \end{align*}
  を満たすことを示す。

  まず、(2)と(3)の式を加える：
  \begin{align*}
    \bbn{J_n}{z}+\frac{n}{z}J_n
    &=J_{n-1}。
  \end{align*}
  したがって：
  \begin{align*}
    J_{n-1}
    =\bbn{J_n}{z}+\frac{n}{z}J_n。
  \end{align*}
  同様に、添字$n$を$n-1$へ置き換える：
  \begin{align*}
    J_{n-2}
    =\bbn{J_{n-1}}{z}+\frac{n-1}{z}J_{n-1}。
  \end{align*}
  よって：
  \begin{align*}
    \bbn{J_{n-1}}{z}
    =J_{n-2}-\frac{n-1}{z}J_{n-1}。
  \end{align*}
  また、(3)の漸化式で$n$を$n-1$へ置き換える：
  \begin{align*}
    \frac{2(n-1)}{z}J_{n-1}=J_{n-2}+J_n。
  \end{align*}
  これを$J_{n-2}$について解く：
  \begin{align*}
    J_{n-2}
    =\frac{2(n-1)}{z}J_{n-1}-J_n。
  \end{align*}
  したがって：
  \begin{align*}
    \bbn{J_{n-1}}{z}
    &=\frac{n-1}{z}J_{n-1}-J_n。
  \end{align*}

  次に、
  \begin{align*}
    \bbn{J_n}{z}
    =J_{n-1}-\frac{n}{z}J_n
  \end{align*}
  をもう一度$z$で微分する：
  \begin{align*}
    \nbbn{2}{J_n}{z}
    &=\bbn{J_{n-1}}{z}
    -\frac{n}{z}\bbn{J_n}{z}
    +\frac{n}{z^2}J_n\\
    &=\frac{n-1}{z}J_{n-1}-J_n
    -\frac{n}{z}\bbn{J_n}{z}
    +\frac{n}{z^2}J_n。
  \end{align*}
  ここに
  \begin{align*}
    J_{n-1}
    =\bbn{J_n}{z}+\frac{n}{z}J_n
  \end{align*}
  を代入する：
  \begin{align*}
    \nbbn{2}{J_n}{z}
    &=
    \frac{n-1}{z}
    \left(\bbn{J_n}{z}+\frac{n}{z}J_n\right)
    -J_n
    -\frac{n}{z}\bbn{J_n}{z}
    +\frac{n}{z^2}J_n\\
    &=-\frac{1}{z}\bbn{J_n}{z}
    +\left(\frac{n^2}{z^2}-1\right)J_n。
  \end{align*}
  すべての項を左辺へ移項する：
  \begin{align*}
    \nbbn{2}{J_n}{z}
    +\frac{1}{z}\bbn{J_n}{z}
    +\left(1-\frac{n^2}{z^2}\right)J_n
    =0。
  \end{align*}
  以上より、(2)および(3)の漸化式を用いることにより、$J_n(z)$が$n$次のBesselの微分方程式を満たすことが示された。

  \subsection*{問題1.10}

  この問題では、矩形関数
  \begin{align*}
    f(x)=
    \begin{cases}
      1,& |x|<1,\\
      0,& |x|>1
    \end{cases}
  \end{align*}
  の自分自身とのたたみ込み $(f*f)(x)$ を求め、さらに Fourier 変換 $F(k)$ を計算し、たたみ込み定理が実際に成り立つことを確かめる。
  \begin{align*}
    \mathcal{F}[f*f](k)=F(k)^2。
  \end{align*}
  まず、たたみ込みの定義より、
  \begin{align*}
    (f*f)(x)
    =
    \int_{-\infty}^{\infty}
    f(x-y)f(y)\,dy。
  \end{align*}
  ここで被積分関数 $f(x-y)f(y)$ は、
  \begin{align*}
    |y|<1,
    \qquad
    |x-y|<1
  \end{align*}
  を同時に満たすときのみ $1$ となり、それ以外では $0$ となる。したがって $(f*f)(x)$ は、区間$(-1,1)$と区間$(x-1,x+1)$の共通部分の長さに等しい。\\
  まず、$|x|\ge 2$の場合を考える。このとき2つの区間は互いに重ならないので、
  \begin{align*}
    (f*f)(x)=0。
  \end{align*}
  次に$0\le x<2$の場合を考える。このとき共通部分は
  \begin{align*}
    x-1<y<1
  \end{align*}
  であり、その長さは
  \begin{align*}
    1-(x-1)
    =
    2-x。
  \end{align*}
  同様に$-2<x<0$の場合には共通部分は
  \begin{align*}
    -1<y<x+1
  \end{align*}
  となるので、その長さは、
  \begin{align*}
    (x+1)-(-1)
    =
    x+2。
  \end{align*}
  以上をまとめると、
  \begin{align*}
    (f*f)(x)
    =
    \begin{cases}
      x+2,& -2<x<0,\\
      2-x,& 0\le x<2,\\
      0,& |x|\ge 2
    \end{cases}
  \end{align*}
  を得る。さらに、この結果は絶対値を用いて
  \begin{align*}
    (f*f)(x)
    =
    \begin{cases}
      2-|x|,& |x|<2,\\
      0,& |x|\ge 2
    \end{cases}
  \end{align*}
  と書くことができる。すなわち、矩形関数を自分自身とたたみ込むと三角形状の関数が得られる。\\
  次にFourier変換を求める。Fourier変換の定義
  \begin{align*}
    F(k)
    =
    \int_{-\infty}^{\infty}
    f(x)e^{-ikx}\,dx
  \end{align*}
  を用いると、
  \begin{align*}
    F(k)
    &=
    \int_{-1}^{1}
    e^{-ikx}\,dx \\
    &=
    \left[
      \frac{e^{-ikx}}{-ik}
    \right]_{-1}^{1} \\
    &=
    \frac{e^{ik}-e^{-ik}}{ik}。
  \end{align*}
  ここでEulerの公式
  \begin{align*}
    e^{ik}-e^{-ik}
    =
    2i\sin k
  \end{align*}
  を用いると、
  \begin{align*}
    F(k)
    =
    \frac{2\sin k}{k}
  \end{align*}
  を得る。したがって、
  \begin{align*}
    F(k)^2
    =
    \frac{4\sin^2 k}{k^2}。
  \end{align*}
  一方、求めたたたみ込み
  \begin{align*}
    (f*f)(x)
    =
    \begin{cases}
      2-|x|,& |x|<2,\\
      0,& |x|\ge 2
    \end{cases}
  \end{align*}
  のFourier変換を直接計算する。\\
  関数は偶関数であるから、以下のように書ける：
  \begin{align*}
    \mathcal{F}[f*f](k)
    &=
    \int_{-2}^{2}
    (2-|x|)e^{-ikx}\,dx \\
    &=
    2\int_0^2
    (2-x)\cos(kx)\,dx。
  \end{align*}
  部分積分を用いて積分を計算する：
  \begin{align*}
    \int_0^2
    (2-x)\cos(kx)\,dx
    &=
    \left[
      (2-x)\frac{\sin(kx)}{k}
    \right]_0^2
    +
    \frac{1}{k}
    \int_0^2
    \sin(kx)\,dx \\
    &=
    \frac{1}{k}
    \left[
      -\frac{\cos(kx)}{k}
    \right]_0^2 \\
    &=
    \frac{1-\cos(2k)}{k^2}。
  \end{align*}
  よって、
  \begin{align*}
    \mathcal{F}[f*f](k)
    &=
    \frac{2\{1-\cos(2k)\}}{k^2} \\
    &=
    \frac{4\sin^2 k}{k^2}。
  \end{align*}
  ここで三角恒等式を用いた：
  \begin{align*}
    1-\cos(2k)
    =
    2\sin^2 k。
  \end{align*}
  したがって、
  \begin{align*}
    \mathcal{F}[f*f](k)
    =
    \frac{4\sin^2 k}{k^2}
    =
    F(k)^2
  \end{align*}
  が確かに成り立つ。\\
  以上より、
  \begin{align*}
    (f*f)(x)
    =
    \begin{cases}
      2-|x|,& |x|<2,\\
      0,& |x|\ge 2
    \end{cases}
  \end{align*}
  であり、そのFourier変換は、
  \begin{align*}
    \mathcal{F}[f*f](k)
    =
    \frac{4\sin^2 k}{k^2}
    =
    F(k)^2
  \end{align*}
  となる。

  \subsection*{問題1.12(a)}

  この小問では、周期$2\pi$をもつ関数
  \begin{align*}
    \delta_n(x)
    =
    \begin{cases}
      n & \left(|x|<\dfrac{1}{2n}\right),\\
      0 & \left(|x|>\dfrac{1}{2n}\right)
    \end{cases}
  \end{align*}
  のFourier級数展開を求める。ただし、$|x|=1/(2n)$における関数値はFourier係数に影響しないため、任意に定めてよい。
  
  まず、$\delta_n(x)$は次の関係を満たす：
  \begin{align*}
    \delta_n(-x)=\delta_n(x)。
  \end{align*}
  したがって、$\delta_n(x)$は偶関数である。一般の周期$2\pi$のFourier級数を
  \begin{align*}
    \delta_n(x)=\frac{a_0}{2}+\sum_{m=1}^{\infty}a_m\cos(mx)+\sum_{m=1}^{\infty}b_m\sin(mx)
  \end{align*}
  と書くと、$\delta_n(x)\sin(mx)$は奇関数であるため、正弦項の係数は次のようになる：
  \begin{align*}
    b_m
    &=\frac{1}{\pi}\int_{-\pi}^{\pi}\delta_n(x)\sin(mx)\,dx\\
    &=0\quad(m=1,2,\dots)。
  \end{align*}
  次に、定数項の係数$a_0$を求める。Fourier係数の定義を用いると、$\delta_n(x)$が非零となる区間は$-1/(2n)<x<1/(2n)$であるため：
  \begin{align*}
    a_0
    &=\frac{1}{\pi}\int_{-\pi}^{\pi}\delta_n(x)\,dx\\
    &=\frac{n}{\pi}\int_{-\frac{1}{2n}}^{\frac{1}{2n}}dx\\
    &=\frac{n}{\pi}\left\{\frac{1}{2n}-\left(-\frac{1}{2n}\right)\right\}\\
    &=\frac{1}{\pi}。
  \end{align*}
  続いて、$m\geq 1$に対する余弦項の係数$a_m$を求める。Fourier係数の定義より：
  \begin{align*}
    a_m
    &=\frac{1}{\pi}\int_{-\pi}^{\pi}\delta_n(x)\cos(mx)\,dx\\
    &=\frac{n}{\pi}\int_{-\frac{1}{2n}}^{\frac{1}{2n}}\cos(mx)\,dx\\
    &=\frac{n}{\pi}\left[\frac{\sin(mx)}{m}\right]_{-\frac{1}{2n}}^{\frac{1}{2n}}\\
    &=\frac{n}{\pi m}\left\{\sin\left(\frac{m}{2n}\right)-\sin\left(-\frac{m}{2n}\right)\right\}。
  \end{align*}
  ここで、正弦関数の奇関数性$\sin(-\theta)=-\sin\theta$を用いる：
  \begin{align*}
    a_m=\frac{2n}{\pi m}\sin\left(\frac{m}{2n}\right)\quad(m=1,2,\dots)。
  \end{align*}
  以上のFourier係数を代入することにより、$\delta_n(x)$のFourier級数を得る：
  \begin{align*}
    \delta_n(x)
    =\frac{1}{2\pi}
    +\sum_{m=1}^{\infty}\frac{2n}{\pi m}\sin\left(\frac{m}{2n}\right)\cos(mx)。
  \end{align*}
  この級数は、$\delta_n(x)$が連続な点では$\delta_n(x)$に収束する。一方、不連続点$x=\pm1/(2n)$およびその周期的な移動で得られる点では、左右極限の平均値
  \begin{align*}
    \frac{n+0}{2}=\frac{n}{2}
  \end{align*}
  に収束する。以上より、$\delta_n(x)$のFourier級数展開が示された。

  \subsection*{問題1.12(b)}

  この小問では、(a)で求めたFourier級数を用いて、$\delta_n(x)$が$n\longrightarrow\infty$の極限において周期的デルタ関数$\delta_P(x)$へ超関数の意味で収束することを示す。
  
  まず、通常の関数としての各点収束を確認しておく。$x\notin2\pi\mathbb{Z}$を固定すると、十分大きな$n$に対して$|x-2\pi k|>1/(2n)$がすべての$k\in\mathbb{Z}$について成り立つため：
  \begin{align*}
    \lim_{n\to\infty}\delta_n(x)=0\quad(x\notin2\pi\mathbb{Z})。
  \end{align*}
  一方、$x\in2\pi\mathbb{Z}$では$\delta_n(x)=n$であるため、その値は発散する。したがって、$\delta_n(x)$は通常の関数として点ごとに$\delta_P(x)$へ収束するわけではない。以下では、周期的な試験関数に作用させた積分の極限によって、超関数としての収束を示す。\\
  (a)より、$\delta_n(x)$のFourier級数は次のように表される：
  \begin{align*}
    \delta_n(x)
    =\frac{1}{2\pi}
    +\sum_{m=1}^{\infty}\frac{2n}{\pi m}\sin\left(\frac{m}{2n}\right)\cos(mx)。
  \end{align*}
  ここで、$m\geq1$を固定して$n\longrightarrow\infty$の極限を考える。正弦関数の原点まわりのTaylor展開を用いる：
  \begin{align*}
    \sin y=y+O(y^3)\quad(y\longrightarrow0)。
  \end{align*}
  ここに$y=m/(2n)$を代入する。$m$を固定しているため：
  \begin{align*}
    \sin\left(\frac{m}{2n}\right)
    =\frac{m}{2n}+O\left(\frac{m^3}{n^3}\right)。
  \end{align*}
  したがって、各Fourier係数について：
  \begin{align*}
    \frac{2n}{\pi m}\sin\left(\frac{m}{2n}\right)
    &=\frac{2n}{\pi m}\left\{\frac{m}{2n}+O\left(\frac{m^3}{n^3}\right)\right\}\\
    &=\frac{1}{\pi}+O\left(\frac{m^2}{n^2}\right)。
  \end{align*}
  よって、$m$を固定したとき：
  \begin{align*}
    \lim_{n\to\infty}\frac{2n}{\pi m}\sin\left(\frac{m}{2n}\right)=\frac{1}{\pi}。
  \end{align*}
  したがって、$\delta_n$の各Fourier係数は、形式的には次のFourier級数の係数へ収束する：
  \begin{align*}
    \frac{1}{2\pi}+\frac{1}{\pi}\sum_{m=1}^{\infty}\cos(mx)
    =\frac{1}{2\pi}\left\{1+2\sum_{m=1}^{\infty}\cos(mx)\right\}。
  \end{align*}
  Eulerの公式より：
  \begin{align*}
    e^{imx}+e^{-imx}=2\cos(mx)
  \end{align*}
  が成り立つため、この形式的な級数は次のように書き換えられる：
  \begin{align*}
    \frac{1}{2\pi}\left\{1+2\sum_{m=1}^{\infty}\cos(mx)\right\}
    =\frac{1}{2\pi}\sum_{m=-\infty}^{\infty}e^{imx}。
  \end{align*}
  右辺は周期的デルタ関数$\delta_P(x)$のFourier級数表示である：
  \begin{align*}
    \delta_P(x)=\frac{1}{2\pi}\sum_{m=-\infty}^{\infty}e^{imx}。
  \end{align*}
  ただし、以上の係数ごとの極限だけでは、無限和と$n\longrightarrow\infty$の極限を交換できることまでは直ちに従わない。そこで、周期的デルタ関数への収束を、周期$2\pi$の連続な試験関数$\varphi(x)$に対する作用によって直接確認する。\\
  $\delta_n(x)$が区間$[-1/(2n),1/(2n)]$でのみ非零であることを用いる：
  \begin{align*}
    \int_{-\pi}^{\pi}\delta_n(x)\varphi(x)\,dx
    =n\int_{-\frac{1}{2n}}^{\frac{1}{2n}}\varphi(x)\,dx。
  \end{align*}
  ここで、$u=nx$と変数変換する。このとき、$x=u/n$および$dx=du/n$であるため：
  \begin{align*}
    \int_{-\pi}^{\pi}\delta_n(x)\varphi(x)\,dx
    &=\int_{-\frac{1}{2}}^{\frac{1}{2}}\varphi\left(\frac{u}{n}\right)\,du。
  \end{align*}
  $\varphi$は連続であるため、各$u\in[-1/2,1/2]$について：
  \begin{align*}
    \lim_{n\to\infty}\varphi\left(\frac{u}{n}\right)=\varphi(0)
  \end{align*}
  が成り立つ。また、$\varphi$はコンパクトな区間$[-\pi,\pi]$上で有界であるから、Lebesgueの優収束定理を適用できる。したがって：
  \begin{align*}
    \lim_{n\to\infty}\int_{-\pi}^{\pi}\delta_n(x)\varphi(x)\,dx
    &=\int_{-\frac{1}{2}}^{\frac{1}{2}}\lim_{n\to\infty}\varphi\left(\frac{u}{n}\right)\,du\\
    &=\int_{-\frac{1}{2}}^{\frac{1}{2}}\varphi(0)\,du\\
    &=\varphi(0)。
  \end{align*}
  一方、周期的デルタ関数の定義より：
  \begin{align*}
    \int_{-\pi}^{\pi}\delta_P(x)\varphi(x)\,dx=\varphi(0)。
  \end{align*}
  よって、任意の周期$2\pi$の連続な試験関数$\varphi$に対して：
  \begin{align*}
    \lim_{n\to\infty}\int_{-\pi}^{\pi}\delta_n(x)\varphi(x)\,dx
    =\int_{-\pi}^{\pi}\delta_P(x)\varphi(x)\,dx。
  \end{align*}
  これは、$\delta_n$が周期的超関数の意味で$\delta_P$へ収束することを表す。したがって：
  \begin{align*}
    \delta_n\xrightarrow[n\to\infty]{}\delta_P
    \quad\text{（周期的超関数の意味で）}。
  \end{align*}
  以上より、(a)で得たFourier係数は周期的デルタ関数のFourier係数へ項別に収束し、さらに$\delta_n$そのものも超関数の意味で周期的デルタ関数$\delta_P$へ収束することが示された。

  \subsection*{問題2.5}

  この問題では、3次元波動方程式の遅延Green関数を$z$方向について積分することにより、2次元波動方程式の遅延Green関数
  \begin{align*}
    G(x,y,t;x',y',t')
    =
    \frac{c\,\theta\!\left(c(t-t')-\sqrt{(x-x')^2+(y-y')^2}\right)}
    {2\pi\sqrt{c^2(t-t')^2-(x-x')^2-(y-y')^2}}
  \end{align*}
  を導く。
  
  3次元波動方程式の遅延Green関数は次のように与えられる：
  \begin{align*}
    G_3(\vect{r},t;\vect{r}\,',t')
    =
    \frac{1}{4\pi|\vect{r}-\vect{r}\,'|}
    \delta\left(t-t'-\frac{|\vect{r}-\vect{r}\,'|}{c}\right)。
  \end{align*}
  ここで、位置ベクトルを次のように置く：
  \begin{align*}
    \vect{r}=(x,y,z),
    \qquad
    \vect{r}\,'=(x',y',z')。
  \end{align*}
  2次元Green関数は、3次元Green関数を$z$方向について積分することにより得られる：
  \begin{align*}
    G(x,y,t;x',y',t')
    =
    \int_{-\infty}^{\infty}G_3(\vect{r},t;\vect{r}\,',t')\,dz。
  \end{align*}
  以下では、記述を簡潔にするため、平面内の2点間の距離$\rho$と時間差$\tau$を次のように定義する：
  \begin{align*}
    \rho=\sqrt{(x-x')^2+(y-y')^2},
    \qquad
    \tau=t-t'。
  \end{align*}
  さらに、$z$方向の相対座標を
  \begin{align*}
    \xi=z-z'
  \end{align*}
  と置く。このとき、$dz=d\xi$であり、3次元空間における2点間の距離は次のように表される：
  \begin{align*}
    |\vect{r}-\vect{r}\,'|=\sqrt{\rho^2+\xi^2}。
  \end{align*}
  したがって、2次元Green関数は次のように書き換えられる：
  \begin{align*}
    G
    =
    \frac{1}{4\pi}
    \int_{-\infty}^{\infty}
    \frac{\delta\left(\tau-\dfrac{\sqrt{\rho^2+\xi^2}}{c}\right)}
    {\sqrt{\rho^2+\xi^2}}\,d\xi。
  \end{align*}
  この積分を計算するため、単純零点$\xi_i$をもつ関数$g(\xi)$について成り立つデルタ関数の公式
  \begin{align*}
    \delta(g(\xi))
    =
    \sum_i\frac{\delta(\xi-\xi_i)}{|g'(\xi_i)|},
    \qquad
    g(\xi_i)=0
  \end{align*}
  を用いる。いま、デルタ関数の引数を
  \begin{align*}
    g(\xi)=\tau-\frac{\sqrt{\rho^2+\xi^2}}{c}
  \end{align*}
  と置く。\\
  まず、$g(\xi)=0$を満たす零点を求める。零点では次の関係が成り立つ：
  \begin{align*}
    \tau
    &=\frac{\sqrt{\rho^2+\xi^2}}{c},\\
    c^2\tau^2
    &=\rho^2+\xi^2。
  \end{align*}
  したがって、零点の候補は次の2点である：
  \begin{align*}
    \xi_\pm=\pm\sqrt{c^2\tau^2-\rho^2}。
  \end{align*}
  右辺の平方根が実数となるためには$c^2\tau^2\geq\rho^2$が必要である。さらに、零点の方程式の右辺は非負であるため、$\tau<0$の場合には実数解が存在しない。したがって、異なる2つの単純零点が存在する領域は
  \begin{align*}
    \tau>0,
    \qquad
    c\tau>\rho
  \end{align*}
  である。\\
  次に、$g(\xi)$を$\xi$で微分する：
  \begin{align*}
    g'(\xi)
    =-\frac{\xi}{c\sqrt{\rho^2+\xi^2}}。
  \end{align*}
  零点$\xi_\pm$においては、
  \begin{align*}
    \sqrt{\rho^2+\xi_\pm^2}=c\tau,
    \qquad
    |\xi_\pm|=\sqrt{c^2\tau^2-\rho^2}
  \end{align*}
  が成り立つ。したがって、零点における導関数の絶対値は次のようになる：
  \begin{align*}
    |g'(\xi_\pm)|
    &=\frac{|\xi_\pm|}{c\sqrt{\rho^2+\xi_\pm^2}}\\
    &=\frac{\sqrt{c^2\tau^2-\rho^2}}{c^2\tau}。
  \end{align*}
  よって、$c\tau>\rho$の領域では、デルタ関数の変数変換公式より：
  \begin{align*}
    \delta(g(\xi))
    =
    \frac{c^2\tau}{\sqrt{c^2\tau^2-\rho^2}}
    \left\{
      \delta\left(\xi-\sqrt{c^2\tau^2-\rho^2}\right)
      +
      \delta\left(\xi+\sqrt{c^2\tau^2-\rho^2}\right)
    \right\}。
  \end{align*}
  これをGreen関数の積分表示へ代入する。被積分関数であるデルタ関数において、引数が非零となる$\xi$に対しては、$\sqrt{\rho^2+\xi^2}$をどのようにおいても被積分関数はゼロのままであるため、積分結果は変わらない。ゆえに、
  \begin{align*}
    \sqrt{\rho^2+\xi^2}=c\tau
  \end{align*}
  と固定して良い。よって、
  \begin{align*}
    G
    &=
    \frac{1}{4\pi}
    \int_{-\infty}^{\infty}
    \frac{1}{c\tau}
    \frac{c^2\tau}{\sqrt{c^2\tau^2-\rho^2}}
    \left\{
      \delta\left(\xi-\sqrt{c^2\tau^2-\rho^2}\right)
      +
      \delta\left(\xi+\sqrt{c^2\tau^2-\rho^2}\right)
    \right\}d\xi\\
    &=
    \frac{c}{4\pi\sqrt{c^2\tau^2-\rho^2}}
    \int_{-\infty}^{\infty}
    \left\{
      \delta\left(\xi-\sqrt{c^2\tau^2-\rho^2}\right)
      +
      \delta\left(\xi+\sqrt{c^2\tau^2-\rho^2}\right)
    \right\}d\xi\\
    &=
    \frac{c}{2\pi\sqrt{c^2\tau^2-\rho^2}}
    \qquad(c\tau>\rho)。
  \end{align*}
  一方、$c\tau<\rho$ではデルタ関数の引数が零となる実数$\xi$が存在しないため、Green関数はゼロである。以上の結果をHeaviside関数によってまとめる：
  \begin{align*}
    G
    =
    \frac{c\,\theta(c\tau-\rho)}
    {2\pi\sqrt{c^2\tau^2-\rho^2}}。
  \end{align*}
  ここで、$\rho\geq0$であるため、$\theta(c\tau-\rho)$は$c\tau\geq\rho$と同時に$\tau\geq0$も保証している。したがって、時間方向の因果条件を表す$\theta(\tau)$を別に掛ける必要はない。また、波面$c\tau=\rho$では分母が発散するが、この特異性は局所可積分であり、上式は分布の意味で解釈される。\\
  最後に、
  \begin{align*}
    \tau=t-t',
    \qquad
    \rho=\sqrt{(x-x')^2+(y-y')^2}
  \end{align*}
  を代入する：
  \begin{align*}
    G(x,y,t;x',y',t')
    =
    \frac{c\,\theta\left(c(t-t')-\sqrt{(x-x')^2+(y-y')^2}\right)}
    {2\pi\sqrt{c^2(t-t')^2-(x-x')^2-(y-y')^2}}。
  \end{align*}
  以上より、3次元波動方程式の遅延Green関数を$z$方向について積分することにより、式(2.196)で与えられる2次元波動方程式の遅延Green関数が導かれた。

  \subsection*{問題5.1}

  この問題では、Legendre多項式$P_n(x)$の生成関数を用いて、$x=-1$における値$P_n(-1)$を求める。
  
  Legendre多項式の生成関数は、$|t|<1$において次のように与えられる：
  \begin{align*}
    G(x,t)
    =\frac{1}{\sqrt{1-2xt+t^2}}
    =\sum_{n=0}^{\infty}P_n(x)t^n。
  \end{align*}
  ここに$x=-1$を代入する：
  \begin{align*}
    G(-1,t)
    &=\frac{1}{\sqrt{1+2t+t^2}}\\
    &=\frac{1}{\sqrt{(1+t)^2}}。
  \end{align*}
  一般に、実数$t$に対して$\sqrt{(1+t)^2}=|1+t|$である。ここでは生成関数を$|t|<1$の範囲で考えているため、$1+t>0$が成り立つ。したがって：
  \begin{align*}
    \sqrt{(1+t)^2}=1+t。
  \end{align*}
  よって、$x=-1$における生成関数は次のように書ける：
  \begin{align*}
    G(-1,t)=\frac{1}{1+t}。
  \end{align*}
  次に、$|t|<1$における等比級数の公式を用いる：
  \begin{align*}
    \frac{1}{1+t}
    =\frac{1}{1-(-t)}
    =\sum_{n=0}^{\infty}(-t)^n
    =\sum_{n=0}^{\infty}(-1)^nt^n。
  \end{align*}
  一方、Legendre多項式の生成関数の定義に$x=-1$を代入すると：
  \begin{align*}
    G(-1,t)=\sum_{n=0}^{\infty}P_n(-1)t^n。
  \end{align*}
  したがって、$|t|<1$において次の2つの冪級数表示が一致する：
  \begin{align*}
    \sum_{n=0}^{\infty}P_n(-1)t^n
    =\sum_{n=0}^{\infty}(-1)^nt^n。
  \end{align*}
  冪級数の係数の一意性より、各次数の係数は一致しなければならない。したがって：
  \begin{align*}
    P_n(-1)=(-1)^n\quad(n=0,1,2,\dots)。
  \end{align*}
  以上より、Legendre多項式の$x=-1$における値は$P_n(-1)=(-1)^n$であることが示された。

  \subsection*{問題9.1}

  この問題では、Laguerre多項式$L_n(x)$を用いて定義した関数
  \begin{align*}
    \phi_n(x)=e^{-x/2}L_n(x)
  \end{align*}
  が、次の微分方程式を満たすことを示す：
  \begin{align*}
    x\nbbn{2}{\phi_n}{x}
    +\bbn{\phi_n}{x}
    +\left(n+\frac{1}{2}-\frac{x}{4}\right)\phi_n
    =0。
  \end{align*}
  まず、Laguerre多項式$L_n(x)$は、Laguerreの微分方程式
  \begin{align*}
    x\nbbn{2}{L_n}{x}
    +(1-x)\bbn{L_n}{x}
    +nL_n
    =0
  \end{align*}
  を満たす。ここで、$\phi_n(x)=e^{-x/2}L_n(x)$を$x$について微分する。積の微分を用いると、1階導関数は次のようになる：
  \begin{align*}
    \bbn{\phi_n}{x}
    &=-\frac{1}{2}e^{-x/2}L_n+e^{-x/2}\bbn{L_n}{x}\\
    &=e^{-x/2}\left(\bbn{L_n}{x}-\frac{1}{2}L_n\right)。
  \end{align*}
  さらに、上式をもう一度$x$について微分する：
  \begin{align*}
    \nbbn{2}{\phi_n}{x}
    &=-\frac{1}{2}e^{-x/2}\left(\bbn{L_n}{x}-\frac{1}{2}L_n\right)
    +e^{-x/2}\left(\nbbn{2}{L_n}{x}-\frac{1}{2}\bbn{L_n}{x}\right)\\
    &=e^{-x/2}\left(\nbbn{2}{L_n}{x}-\bbn{L_n}{x}+\frac{1}{4}L_n\right)。
  \end{align*}
  以上で得られた$\phi_n$の1階導関数および2階導関数を、示すべき微分方程式の左辺へ代入する：
  \begin{align*}
    &x\nbbn{2}{\phi_n}{x}
    +\bbn{\phi_n}{x}
    +\left(n+\frac{1}{2}-\frac{x}{4}\right)\phi_n\\
    &\quad
    =xe^{-x/2}\left(\nbbn{2}{L_n}{x}-\bbn{L_n}{x}+\frac{1}{4}L_n\right)
    +e^{-x/2}\left(\bbn{L_n}{x}-\frac{1}{2}L_n\right)
    +e^{-x/2}\left(n+\frac{1}{2}-\frac{x}{4}\right)L_n\\
    &\quad
    =e^{-x/2}\left\{
      x\nbbn{2}{L_n}{x}
      -x\bbn{L_n}{x}
      +\frac{x}{4}L_n
      +\bbn{L_n}{x}
      -\frac{1}{2}L_n
      +nL_n
      +\frac{1}{2}L_n
      -\frac{x}{4}L_n
    \right\}\\
    &\quad
    =e^{-x/2}\left\{
      x\nbbn{2}{L_n}{x}
      +(1-x)\bbn{L_n}{x}
      +nL_n
    \right\}。
  \end{align*}
  最後の中括弧内はLaguerreの微分方程式の左辺に一致するため：
  \begin{align*}
    x\nbbn{2}{L_n}{x}
    +(1-x)\bbn{L_n}{x}
    +nL_n
    =0。
  \end{align*}
  したがって、次の式を得る：
  \begin{align*}
    x\nbbn{2}{\phi_n}{x}
    +\bbn{\phi_n}{x}
    +\left(n+\frac{1}{2}-\frac{x}{4}\right)\phi_n
    =0。
  \end{align*}
  以上より、関数$\phi_n(x)=e^{-x/2}L_n(x)$が、与えられた微分方程式を満たすことが示された。

  \subsection*{問題10.2}

  この問題では、既知の半整数次Bessel関数と漸化式を用いて、$J_{5/2}(z)$および$J_{-5/2}(z)$を求める。以下では、平方根を実数として扱うため、$z>0$とする。
  
  まず、$\pm 1/2$次Bessel関数は次のように与えられる：
  \begin{align*}
    J_{1/2}(z)
    &=\sqrt{\frac{2}{\pi z}}\sin z,\\
    J_{-1/2}(z)
    &=\sqrt{\frac{2}{\pi z}}\cos z。
  \end{align*}
  また、Bessel関数の漸化式を用いる：
  \begin{align*}
    J_{\nu+1}(z)=\frac{2\nu}{z}J_\nu(z)-J_{\nu-1}(z)。
  \end{align*}
  まず、$J_{5/2}(z)$を求める。上の漸化式に$\nu=1/2$を代入する：
  \begin{align*}
    J_{3/2}(z)
    &=\frac{1}{z}J_{1/2}(z)-J_{-1/2}(z)\\
    &=\sqrt{\frac{2}{\pi z}}
    \left(\frac{\sin z}{z}-\cos z\right)。
  \end{align*}
  次に、漸化式に$\nu=3/2$を代入する：
  \begin{align*}
    J_{5/2}(z)
    &=\frac{3}{z}J_{3/2}(z)-J_{1/2}(z)\\
    &=\frac{3}{z}\sqrt{\frac{2}{\pi z}}
    \left(\frac{\sin z}{z}-\cos z\right)
    -\sqrt{\frac{2}{\pi z}}\sin z\\
    &=\sqrt{\frac{2}{\pi z}}
    \left\{\left(\frac{3}{z^2}-1\right)\sin z-\frac{3}{z}\cos z\right\}。
  \end{align*}
  したがって、$J_{5/2}(z)$は次のように得られる：
  \begin{align*}
    J_{5/2}(z)
    =\sqrt{\frac{2}{\pi z}}
    \left\{\left(\frac{3}{z^2}-1\right)\sin z-\frac{3}{z}\cos z\right\}。
  \end{align*}
  次に、$J_{-5/2}(z)$を求める。このため、漸化式を次のように書き換える：
  \begin{align*}
    J_{\nu-1}(z)=\frac{2\nu}{z}J_\nu(z)-J_{\nu+1}(z)。
  \end{align*}
  この式に$\nu=-1/2$を代入する：
  \begin{align*}
    J_{-3/2}(z)
    &=-\frac{1}{z}J_{-1/2}(z)-J_{1/2}(z)\\
    &=-\sqrt{\frac{2}{\pi z}}
    \left(\frac{\cos z}{z}+\sin z\right)。
  \end{align*}
  さらに、同じ漸化式に$\nu=-3/2$を代入する：
  \begin{align*}
    J_{-5/2}(z)
    &=-\frac{3}{z}J_{-3/2}(z)-J_{-1/2}(z)\\
    &=\frac{3}{z}\sqrt{\frac{2}{\pi z}}
    \left(\frac{\cos z}{z}+\sin z\right)
    -\sqrt{\frac{2}{\pi z}}\cos z\\
    &=\sqrt{\frac{2}{\pi z}}
    \left\{\frac{3}{z}\sin z+\left(\frac{3}{z^2}-1\right)\cos z\right\}。
  \end{align*}
  したがって、$J_{-5/2}(z)$は次のように得られる：
  \begin{align*}
    J_{-5/2}(z)
    =\sqrt{\frac{2}{\pi z}}
    \left\{\frac{3}{z}\sin z+\left(\frac{3}{z^2}-1\right)\cos z\right\}。
  \end{align*}
  以上より、求める半整数次Bessel関数は次のようになる：
  \begin{align*}
    J_{5/2}(z)
    &=\sqrt{\frac{2}{\pi z}}
    \left\{\left(\frac{3}{z^2}-1\right)\sin z-\frac{3}{z}\cos z\right\},\\
    J_{-5/2}(z)
    &=\sqrt{\frac{2}{\pi z}}
    \left\{\frac{3}{z}\sin z+\left(\frac{3}{z^2}-1\right)\cos z\right\}。
  \end{align*}
  以上より、$J_{5/2}(z)$および$J_{-5/2}(z)$が求められた。

  \subsection*{問題11.1}

  この問題では、双曲線関数$\cosh(kx)$および$\sinh(kx)$のLaplace変換を求める。関数$f(x)$のLaplace変換は、積分が収束する$s$の範囲において次のように定義される：
  \begin{align*}
    \mathcal{L}\{f\}(s)=\int_0^\infty e^{-sx}f(x)\,dx。
  \end{align*}
  
  まず、$\cosh(kx)$のLaplace変換を求める。双曲線余弦関数の定義
  \begin{align*}
    \cosh(kx)=\frac{e^{kx}+e^{-kx}}{2}
  \end{align*}
  を用いる：
  \begin{align*}
    \mathcal{L}\{\cosh(kx)\}(s)
    &=\int_0^\infty e^{-sx}\cosh(kx)\,dx\\
    &=\frac{1}{2}\int_0^\infty e^{-sx}\left(e^{kx}+e^{-kx}\right)\,dx\\
    &=\frac{1}{2}\int_0^\infty\left\{e^{-(s-k)x}+e^{-(s+k)x}\right\}\,dx。
  \end{align*}
  ここで、2つの積分がともに収束するためには：
  \begin{align*}
    s-k>0,
    \qquad
    s+k>0
  \end{align*}
  が必要である。したがって、$s,k\in\mathbb{R}$の場合の収束条件は次のようになる：
  \begin{align*}
    s>|k|。
  \end{align*}
  この条件のもとでは、各指数関数は$x\longrightarrow\infty$でゼロへ収束する。したがって、積分を実行する：
  \begin{align*}
    \mathcal{L}\{\cosh(kx)\}(s)
    &=\frac{1}{2}\left(\frac{1}{s-k}+\frac{1}{s+k}\right)\\
    &=\frac{1}{2}\frac{(s+k)+(s-k)}{(s-k)(s+k)}\\
    &=\frac{s}{s^2-k^2}。
  \end{align*}
  よって、$\cosh(kx)$のLaplace変換は次のように与えられる：
  \begin{align*}
    \mathcal{L}\{\cosh(kx)\}(s)
    =\frac{s}{s^2-k^2}
    \qquad(s>|k|)。
  \end{align*}
  次に、$\sinh(kx)$のLaplace変換を求める。双曲線正弦関数の定義
  \begin{align*}
    \sinh(kx)=\frac{e^{kx}-e^{-kx}}{2}
  \end{align*}
  を用いる：
  \begin{align*}
    \mathcal{L}\{\sinh(kx)\}(s)
    &=\int_0^\infty e^{-sx}\sinh(kx)\,dx\\
    &=\frac{1}{2}\int_0^\infty e^{-sx}\left(e^{kx}-e^{-kx}\right)\,dx\\
    &=\frac{1}{2}\int_0^\infty\left\{e^{-(s-k)x}-e^{-(s+k)x}\right\}\,dx。
  \end{align*}
  この場合も、2つの指数関数をそれぞれ積分するための収束条件は$s>|k|$である。この条件のもとで積分を実行する：
  \begin{align*}
    \mathcal{L}\{\sinh(kx)\}(s)
    &=\frac{1}{2}\left(\frac{1}{s-k}-\frac{1}{s+k}\right)\\
    &=\frac{1}{2}\frac{(s+k)-(s-k)}{(s-k)(s+k)}\\
    &=\frac{k}{s^2-k^2}。
  \end{align*}
  よって、$\sinh(kx)$のLaplace変換は次のように与えられる：
  \begin{align*}
    \mathcal{L}\{\sinh(kx)\}(s)
    =\frac{k}{s^2-k^2}
    \qquad(s>|k|)。
  \end{align*}
  以上より、求めるLaplace変換は次のようになる：
  \begin{align*}
    \mathcal{L}\{\cosh(kx)\}(s)
    &=\frac{s}{s^2-k^2},\\
    \mathcal{L}\{\sinh(kx)\}(s)
    &=\frac{k}{s^2-k^2}
    \qquad(s>|k|)。
  \end{align*}
  以上より、$\cosh(kx)$および$\sinh(kx)$のLaplace変換が求められた。




\end{document}
